\subsection{Residue-gap arithmetic}

\begin{theorem}[B\'ezout / modular inverse]
  \label{thm:collatz-bezout}
  \lean{Obsmetrologia.Closure.Collatz.bezout}
  \leanok
  \uses{thm:collatz-nondescent-iff}
  B\'ezout's identity over $\mathbb{Z}$ by the extended Euclidean recursion
  (\texttt{Nat.gcd.induction}): $a x + b y = \gcd(a,b)$ for integer
  coefficients. For coprime $D,M$ this gives a modular inverse
  $M \mid D x - 1$ (\texttt{mod\_inverse\_exists}), hence the least residue
  $r<M$ with $D\cdot r\equiv C\pmod M$ (\texttt{exists\_residue}).
  $\mathsf{Classical}$-free.
\end{theorem}

\begin{theorem}[Residue-gap refinement]
  \label{thm:collatz-residue-gap}
  \lean{Obsmetrologia.Closure.Collatz.residue_characterization}
  \leanok
  \uses{thm:collatz-nondescent-iff, thm:collatz-bezout}
  Since $\gcd(D,M)=1$ ($D$ odd, $M$ a power of two), semigroup membership is a
  pure residue/size condition: there is a least residue $r<M$ (constructed via
  the modular inverse, \cref{thm:collatz-bezout}) with
  $D\cdot r\equiv C\pmod M$ and
  \[
    C\in\langle D,M\rangle \iff D\cdot r\le C,
  \]
  so non-membership is the modular inequality $C<D\cdot r$ (equivalently, under
  the congruence, the carry inequality $M\le D\cdot r-C$). Built from a
  coprime-cancellation lemma and the
  $\mathrm{mod}\leftrightarrow\mathrm{dvd}$ bridges over the local congruence
  $a\bmod M=b\bmod M$; $\mathsf{Classical}$-free. The remaining open part is
  purely the gap conjecture $D\cdot r-C\ge M$ for every nontrivial parent
  barrier (\cref{thm:collatz-gap-holds}).
\end{theorem}

\begin{theorem}[Gap reduces to the residue inequality]
  \label{thm:collatz-residue-gap-reduces}
  \lean{Obsmetrologia.Closure.Collatz.parentBarrierGap_iff_parentBarrierResidueGap}
  \leanok
  \uses{thm:collatz-residue-gap, thm:collatz-gap-holds}
  For a parent barrier, $D=2B-A$ is odd and $M=2B=2^{|p|+1}$ a power of two,
  hence coprime (\texttt{parentD\_coprime\_parentM}); so the residue
  characterization applies pointwise. The sharp residue form
  \texttt{ParentBarrierResidueGap} --- for every nontrivial parent barrier the
  residue satisfies $C<D\cdot r$ --- is equivalent to \texttt{ParentBarrierGap}
  (\cref{thm:collatz-gap-holds}). The semigroup-nonmembership part of the Collatz
  frontier is thus isolated, in proved Lean, to a concrete modular inequality per
  barrier. The separate crossing/no-infinite-barrier input remains recorded in
  \cref{thm:collatz-crossing-holds}.
\end{theorem}
\begin{theorem}[Residue gap is the carry inequality]
  \label{thm:collatz-residue-gap-carry}
  \lean{Obsmetrologia.Closure.Collatz.parentBarrierGap_iff_parentBarrierResidueCarryGap}
  \leanok
  \uses{thm:collatz-residue-gap-reduces}
  Under the congruence \(D\cdot r\equiv C\pmod M\) and \(M>0\), Lean proves
  \(C<D\cdot r\) iff \(M\le D\cdot r-C\)
  (\texttt{residue\_lt\_iff\_modulus\_le\_sub}). Therefore the parent-barrier
  semigroup gap is also exactly \texttt{ParentBarrierResidueCarryGap}: every
  nontrivial parent-barrier residue carries by at least one full modulus
  \(M=2^{|p|+1}\). This is the formal high-\(2\)-adic-digit face of the same
  frontier.
\end{theorem}
\begin{theorem}[Residue gap is first-passage cylinder descent]
  \label{thm:collatz-residue-gap-cylinder-descent}
  \lean{Obsmetrologia.Closure.Collatz.firstPassage_seed_residue}
  \lean{Obsmetrologia.Closure.Collatz.parentBarrier_residue_realizes_firstPassage}
  \lean{Obsmetrologia.Closure.Collatz.parentBarrier_residue_descent_iff_gap}
  \lean{Obsmetrologia.Closure.Collatz.parentBarrierResidueGap_iff_cylinderSeeds_descend}
  \leanok
  \uses{thm:collatz-residue-gap-reduces, thm:collatz-cylinder-seeds-descend}
  The modular inequality form is also equivalent to the dynamical first-passage
  descent skeleton: every nontrivial parent-barrier cylinder representative
  descends at its first-passage word iff \texttt{ParentBarrierResidueGap} holds.
  Lean now proves the bridge at the exact residue used by the audit: any seed
  realizing \(p{+}[E]\) solves the residue equation, the unique inverse residue
  \(r\) realizes that first-passage cylinder, and descent of this same \(r\) is
  equivalent to \(C<D\cdot r\). Thus the residue inequality is not merely
  sufficient for first-passage descent; it is exactly the first-passage descent
  frontier expressed in modular arithmetic.
\end{theorem}

\begin{theorem}[Semigroup certificates are bad cylinder seeds]
  \label{thm:collatz-bad-cylinder-language}
  \lean{Obsmetrologia.Closure.Collatz.badCylinderSeed_of_semigroup_certificate}
  \lean{Obsmetrologia.Closure.Collatz.semigroup_certificate_of_badCylinderSeed}
  \lean{Obsmetrologia.Closure.Collatz.inTwoGenSemigroup_iff_badCylinderSeed}
  \lean{Obsmetrologia.Closure.Collatz.parentBarrierGap_iff_no_badCylinderSeed}
  \leanok
  \uses{thm:collatz-residue-gap-cylinder-descent}
  For a fixed parent barrier, \(C\in\langle D,M\rangle\) is equivalent to the
  existence of a seed \(n<M\) in the first-passage cylinder \(p{+}[E]\) whose
  first-passage iterate does not descend: \(n\le T^{|p|+1}(n)\). Therefore the
  remaining parent-barrier gap is also the emptiness of the bad-cylinder language:
  no nontrivial parent-barrier cylinder contains a non-descending representative.
  The backward-seed layer names the same representative directly: the inverse
  residue is the bounded backward seed for \(p{+}[E]\), so the lock can be read
  as emptiness of bad backward seeds. This is the finite-state/sofic attack face
  of the lock; the emptiness itself is still the frontier.
\end{theorem}

\begin{theorem}[Canonical denominator gap is exactly the same-step margin]
  \label{thm:collatz-canonical-denominator-same-step-margin}
  \lean{Obsmetrologia.Closure.Collatz.canonicalDenominator_update_same}
  \lean{Obsmetrologia.Closure.Collatz.canonicalDenominator_update_up}
  \lean{Obsmetrologia.Closure.Collatz.canonicalDenominator_failure_forces_next_up_margin}
  \lean{Obsmetrologia.Closure.Collatz.canonicalDenominator_failure_forces_next_not_same}
  \lean{Obsmetrologia.Closure.Collatz.canonicalDenominatorBound_step_same_iff_margin}
  \lean{Obsmetrologia.Closure.Collatz.canonicalDenominatorBound_iff_same_step_margins}
  \lean{Obsmetrologia.Closure.Collatz.canonicalDenominator_failure_forces_top_quotient}
  \lean{Obsmetrologia.Closure.Collatz.canonicalDenominatorBound_iff_not_top_quotient}
  \lean{Obsmetrologia.Closure.Collatz.canonicalDenominatorBound_iff_not_last_dyadic_slice}
  \lean{Obsmetrologia.Closure.Collatz.canonicalDenominatorBound_all_iff_not_top_quotient}
  \lean{Obsmetrologia.Closure.Collatz.canonicalDenominatorBound_all_iff_not_last_dyadic_slice}
  \lean{Obsmetrologia.Closure.Collatz.canonicalDenominator_failure_iff_top_remainder_exists}
  \lean{Obsmetrologia.Closure.Collatz.canonicalDenominator_no_top_remainder_iff_not_last_dyadic_slice}
  \lean{Obsmetrologia.Closure.Collatz.canonicalDenominatorBound_all_iff_no_top_remainder_exists}
  \leanok
  \uses{thm:collatz-residue-gap-cylinder-descent}
  In the canonical front-loaded parent barrier, write
  \(D=2^{a+b+1}-3^a\), \(P=2^{a+b+1}\), and \(T=2^{b+1}\). Lean proves the
  exact update formulas: the same-\(b\) transition sends \(D\) to \(3D-P\), and
  the up-step transition sends \(D\) to \(P+3D\). Up-steps propagate the
  denominator lower bound automatically. Therefore the global canonical
  denominator-growth claim is equivalent to the same-step margin family
  \(P+T\le 3D\) at every same-\(b\) transition. Lean also proves that any
  failure at \(a\ge2\) lies strictly on the next up-branch side:
  \(3D<P\), so a failed denominator bound cannot itself be followed by another
  same-\(b\) transition. This is the exact formula for the remaining canonical
  arithmetic aperture. Lean also proves the
  counterexample shape: under the parent-barrier upper bound, any failure must
  put \(3^a/2^{b+1}\) in the top quotient block \(2^a-1\), and conversely the
  denominator bound is exactly the assertion that this top quotient block is
  avoided. Equivalently, \(3^a\) must avoid the final dyadic slice of width
  \(2^{b+1}\) immediately below the parent ceiling \(2^{a+b+1}\). This is not
  yet a proof that the top block is never hit. The two avoidance readings are
  also proved equivalent to the full global canonical denominator-growth target
  for every \(a\ge3\). Finally, failure is equivalent to an explicit positive
  remainder \(\rho\) in the forbidden slice:
  \(3^a=(2^a-1)2^{b+1}+\rho\) with \(0<\rho<2^{b+1}\). Globally, the
  denominator-growth target is exactly the absence of such positive top-slice
  remainders at every canonical barrier.
\end{theorem}

\begin{theorem}[Same-step margins close the canonical semigroup gap]
  \label{thm:collatz-canonical-margins-close-gap}
  \lean{Obsmetrologia.Closure.Collatz.CanonicalSameStepMargins}
  \lean{Obsmetrologia.Closure.Collatz.canonicalWord_not_inTwoGenSemigroup_of_sameStepMargins}
  \leanok
  \uses{thm:collatz-canonical-denominator-same-step-margin}
  The same-step margin family is a named proposition
  (\texttt{CanonicalSameStepMargins}), and Lean proves end-to-end that it
  excludes the two-generator semigroup obstruction on every canonical parent
  barrier \(O^a E^b\) with \(a\ge3\): margins give the global denominator
  growth, the denominator bound forces every canonical first-passage backward
  seed to descend, and descent is the gap. The canonical face of the
  parent-barrier frontier is thereby reduced to this one margin family.
\end{theorem}

\begin{theorem}[Interval primorials embed in Pascal's central column]
  \label{thm:nt-interval-primorial-central-column}
  \lean{Obsmetrologia.NumberTheory.choose_mul_factorial_mul_factorial}
  \lean{Obsmetrologia.NumberTheory.Prime.dvd_factorial}
  \lean{Obsmetrologia.NumberTheory.Prime.not_dvd_factorial}
  \lean{Obsmetrologia.NumberTheory.Prime.dvd_choose_central}
  \lean{Obsmetrologia.NumberTheory.Prime.dvd_of_dvd_prod}
  \lean{Obsmetrologia.NumberTheory.prod_dvd_of_pairwise_primes_dvd}
  \lean{Obsmetrologia.NumberTheory.primes_top_interval_prod_dvd_choose_central}
  \lean{Obsmetrologia.NumberTheory.sum_choose}
  \lean{Obsmetrologia.NumberTheory.choose_central_symm}
  \lean{Obsmetrologia.NumberTheory.choose_central_le_two_pow}
  \lean{Obsmetrologia.NumberTheory.primes_top_interval_prod_le}
  \leanok
  The Pascal-recurrence binomial coefficient agrees with the factorial
  quotient (\(C(n,k)\,k!\,(n-k)! = n!\)); a prime divides \(n!\) exactly when
  \(p\le n\); hence every prime of \((m{+}1, 2m{+}1]\) divides the central
  entry \(C(2m{+}1,m)\), and a pairwise-distinct list of such primes has
  product dividing it. The row sum prices the capture: the binomial theorem
  at the discrete carrier with \(x=y=1\) gives \(\sum_k C(n,k)=2^n\), the two
  central entries of an odd row agree, so \(C(2m{+}1,m)\le 2^{2m}\) and the
  interval's primorial is at most \(2^{2m}\). This is the
  divisibility-plus-size skeleton of the prime-counting estimates
  (Erd\H{o}s primorial bound, \(\mathrm{lcm}(1..n)\) growth) that the
  approximant denominators of a logarithm irrationality measure consume.
\end{theorem}

\begin{theorem}[Canonical same-step margins hold]
  \label{thm:collatz-canonical-margins-hold}
  \lean{Obsmetrologia.Closure.Collatz.canonicalSameStepMargins_below_20000}
  \lean{Obsmetrologia.Closure.Collatz.canonicalWord_not_inTwoGenSemigroup_of_lt_20000}
  \notready
  \uses{thm:collatz-canonical-margins-close-gap, thm:nt-interval-primorial-central-column}
  \emph{(Open frontier --- stated, not proved.)} The margin family itself:
  \(P+T\le 3D\) at every same-\(b\) transition with \(a\ge3\). By the proved
  equivalences this says \(3^a\) never lands in the final dyadic slice of
  width \(2^{b+1}\) below \(2^{a+b+1}\) --- the observed window of the
  \(\times 3\) cascade is never the monochrome block \(2^a-1\). Numerically
  clean through \(a=20000\); the tightest margins are \(a=3\) and \(a=5\).
  This range is now kernel-checked in Lean, not merely audited outside it:
  \texttt{canonicalSameStepMargins\_below\_20000} decides the whole margin
  family for every \(3\le a<20000\) by a single near-linear \texttt{decide}
  pass that carries \(\texttt{canonicalTrailingEven}\,a\) forward one step at
  a time (\(O(N)\) kernel work) rather than recomputing it from scratch at
  each \(a\) (the \(O(N^2)\) cost that capped the first kernel-checked
  instance, \texttt{canonicalSameStepMargins\_below\_3000}, at \(a<3000\));
  \texttt{canonicalWord\_not\_inTwoGenSemigroup\_of\_lt\_20000} chains this
  through to the semigroup gap itself for every canonical parent barrier in
  that range, no \texttt{sorry}/\texttt{axiom}/\texttt{native\_decide}
  involved. The intended proof route is an effective irrationality measure for
  \(\log_2 3\) at strength \(2^{-a}\), far below the polynomial strength of
  the Pad\'e/hypergeometric method: integer approximant sequences for the
  logarithm (three-term recurrences over the constructive
  \texttt{logReal} of \texttt{Calculus.Real.Logarithm}), the interval
  primorial estimates above for the denominators, and the integer gap
  (a nonzero integer is at least one) for nonvanishing. An elementary
  continued-fraction certificate was investigated as a lighter alternative
  route and found infeasible even to partially formalize: the recurrence the
  canonical denominator obeys at each same-\(b\)/up-step transition
  (\(D\mapsto 3D-P\) or \(D\mapsto P+3D\)) is exactly the continuant
  recursion for the continued fraction of \(\log_2 3\), so bounding it
  elementarily is the same quantitative-distance-between-powers-of-\(2\)-and-
  \(3\) problem the irrationality-measure route already targets, not a
  shortcut around it. Closing this node
  closes the canonical face of \cref{thm:collatz-gap-holds}. The general
  face is \emph{not} reducible to the canonical family: at length \(4\) the
  barrier \(OOEO\) has margin ratio \(C/(D\,n)=23/55\), strictly worse than
  canonical \(OOOE\)'s \(19/115\), and per length the worst barriers are the
  orbit record-climbers (seeds \(7\), \(11\), \(127\), \(159\), \(239\),
  \ldots), of no fixed word shape. The general face is confined instead by
  the offset envelope (\cref{thm:collatz-offset-envelope}): per barrier,
  failure requires \(D\cdot n\) at most \(2^{L-a}(3^a-2^a)\), so only
  finitely many small seeds per class remain, and those are dynamical, not
  word-combinatorial, content.
\end{theorem}

\begin{theorem}[The offset envelope confines the gap to small seeds]
  \label{thm:collatz-offset-envelope}
  \lean{Obsmetrologia.Closure.Collatz.oddCount_le_length}
  \lean{Obsmetrologia.Closure.Collatz.trace_C_lower_bound}
  \lean{Obsmetrologia.Closure.Collatz.trace_C_upper_bound}
  \lean{Obsmetrologia.Closure.Collatz.firstPassageBackwardSeeds_descend_of_offset_bound}
  \leanok
  \uses{thm:collatz-residue-gap-cylinder-descent}
  Over all parity words with \(a\) odd steps in \(L\) letters, the affine
  offset satisfies \(3^a - 2^a \le C \le 2^{L-a}(3^a - 2^a)\): the
  front-loaded word minimizes and the back-loaded word maximizes, by
  induction on the trailing letter. Consequently every first-passage
  backward seed \(n\) with
  \(2^{L-a}(3^a-2^a) < D\cdot n\) descends unconditionally. The
  parent-barrier gap frontier is thereby confined, per barrier, to the
  finitely many seeds below the envelope threshold.
\end{theorem}
